\section{Estimates and virtual classes at a joined double cover}\label{app:lowerjoins}
\subsection{The local equation for Theorem~\ref{p10-thm:nodes}}\label{app:double-equation}
Fix the joining set and the branch frames of \S\ref{p10-sec:nodes}.
\subsubsection{The deformation space and the matching equations}

At one selected point the nodal map and its deformation are given by
\begin{equation}\label{p10-eq:node-family}
          zy=q,\qquad w=w_0+(z+p)^2+y .
\end{equation}
When $q=0$, the coordinate on the normalized double is $z+p$ and
its attaching mark has coordinate $p$. Translating $w_0$ to zero and
eliminating $y$ gives the free algebra on $1,z,z^2$ with relation
\begin{equation}\label{p10-eq:node-cubic}
                    z(z+p)^2-wz+q=0 .
\end{equation}
To see versality, use the relative cotangent complex of this
hypersurface over the target \cite[Tag~08SL]{stacks}. Its dual on the
central fiber is multiplication by $3z^2-w$ between two copies of
the local ring. This element is a nonzerodivisor on the two branches.
Its cokernel is $\CC\{z\}/(z^3)$, with basis represented by the
variations of $q,w_0,p$, respectively. An unselected branch contributes
one complex parameter. The first deformation sheaf has finite support
and the complex has no other cohomology, so the fixed-target chart is
smooth of complex dimension $b+2j$ and carries its ordinary fundamental
class. Distinct branch disks supply independent pairs $(p_i,q_i)$.
The stabilizer is the single involution of $Y$ over $C$. In
\eqref{p10-eq:node-family} it acts by $(z,p,q)\mapsto(-z,-p,-q)$ at every
join and fixes $s$.

Put $P_I=\operatorname{coker} h_I^*D$. A kernel section restricts to
zero on the identity and to a multiple of $\xi$ on $Y$. Every such pair
matches because $\xi$ vanishes at the marked points. The normalization
sequence for the operator thus identifies $\ker h_I^*D=K$ and gives
\begin{equation}\label{p10-eq:node-matching}
       0\longrightarrow\bigoplus_{i\in I}N_{x_i}
       \longrightarrow P_I\longrightarrow P_f\longrightarrow0 .
\end{equation}
The smoothing calculation takes place in the individual matching fibers
of this sequence. We give it at one point, suppressing $i$.

\subsubsection{The pullback operator on the fixed target}

For a map $h$ in the finite chart, let $F=h_*\mathcal O$ on $C$. It is
a rank-three bundle by \eqref{p10-eq:node-cubic}, the ordinary quadratic
algebra at the remaining branch points, and the sheet decomposition off
the branch disks. Translate the moving branch centers on smaller disks
and extend the translations by diffeomorphisms supported in larger
disks. Identify the family of bundles smoothly, using the polynomial
frames on the smaller disks. On the intervening collars the sheets are
unramified and the identifications are smooth in the parameters.
Fix $2<a<4$. In these trivializations consider the operator
\[
  \mathcal A_h:W^{1,a}(C,N\otimes F)\longrightarrow
       L^a(C,\Omega_C^{0,1}\otimes N\otimes F)
\]
which acts by $D_t$ after evaluation on each regular sheet. The
complex-linear part is the tensor product Dolbeault operator. In the
polynomial frame the antilinear part applies $B$ to the coefficients
and replaces $1,\bar z,\bar z^2$ by their quadratic interpolants at
the three roots. This determines a measurable matrix $M_{p,q}$ almost
everywhere. At an ordinary simple branch the sole quadratic-sheet
factor in the interpolation of $\bar z$ is $\bar z/z$, which is bounded.

Let $d$ be the diameter of the three roots and $m$ their maximum
modulus. In a second divided difference order the two most distant
roots first and last. The first divided differences of conjugation have
modulus one and those of $\bar z^2$ have modulus at most $2m$. Hence
the respective second differences are bounded by $2/d$ and $4m/d$.
Newton interpolation bounds every entry of $M_{p,q}$ by $C(1+d^{-1})$.
Subtracting the mean of the roots depresses \eqref{p10-eq:node-cubic} to a
monic cubic with linear coefficient $-(w+p^2/3)$, so
$|w+p^2/3|\le C d^2$. It follows that the matrices have uniformly bounded
$L^{a'}$ norm for $a<a'<4$. Away from the limiting branch values they
converge pointwise, hence they converge in $L^a$. Since
$W^{1,a}\hookrightarrow C^0$ is compact, the antilinear part is a
continuous family of compact perturbations of the differential part.

We identify its finite-dimensional reduction with that of the actual
pullback operator. A polynomial section in the kernel evaluates to a
bounded solution on the normalization off the branch points and node
preimages. Removability extends it there, and the common root gives
equal values at the two sides of a node. Evaluation is injective by
interpolation on the open unramified locus. This argument also applies
to solutions with smooth forcing supported in that locus. At the
central map the generator on the two local components is represented
by $z u_1(w)$, where $u_1(w)=\xi(z)/z$ on the double. Oddness and
$\bar\partial\xi(0)=0$ give $\xi(z)=v z+O(|z|^3)$ with derivative
error $O(|z|^2)$. Thus $u_1(0)=v$ and its first weak derivatives in
$w=z^2$ are bounded. The same division on each ordinary branch disk
shows that this is a global coefficient solution. The two central
kernels agree.

The real indices agree by deformation to the complex-linear operators
and the finite projection formula, including on the nodal fiber.
Consequently the central cokernels have equal dimension. Smooth forms supported on unramified disks span both cokernels: a
dual annihilator vanishes on these disks and hence off the nodes
by unique continuation. In the normalized matching presentation,
testing arbitrary endpoint values then kills the matching multipliers.
Inside a finite common span surjecting onto both, a generic tuple of
minimal length projects to a basis of each. Use its span $E$ as a
stabilization. Both stabilized operators are surjective near the
central map and their kernels have rank one, using the continuity of
the pullback reduction on nodal curves
\cite[Appendix~B, (B.6)--(B.9)]{bai}. Evaluation identifies the two
lines and commutes with projection to $E$. Normalize a generator by a
fixed real functional supported on an unramified disk and taking value
one on $\xi$. The two obstruction maps to $E$ then coincide. Enlarging
$E$ by further such forms gives the same identification of their
cokernel classes, since evaluation again identifies the stabilized
kernels of equal dimension.

\subsubsection{The logarithmic term}

Set $\ell=\max(|p|,|q|^{1/3})$ and write $Q=q\log(1/|q|)$ for $q\ne0$,
with $Q=0$ at zero. The only entry of $M_{p,q}$ whose $q$-variation
has nonintegrable first derivative is the coefficient of $z^2$ in the
interpolant of $\bar z$. We compute its variation against a H\"older
continuous test using Euclidean area $dA(w)$.

In the disk $|w|\le A\ell^2$, for fixed large $A$, the roots have
size $O(\ell)$ by rescaling \eqref{p10-eq:node-cubic}. The bound by a
constant times $1+|w+p^2/3|^{-1/2}$ gives an integral $O(\ell^3)$
there for the varying and central matrices. On the complementary
annulus choose $r^2=w$ and set $z=rZ$. The roots in $Z$ are analytic
near $0,1,-1$ in $p/r$ and $q/r^3$, which are uniformly small. Their
interpolation matrix is real analytic in these parameters, with bounded
second derivatives. The coefficient in question is multiplied by
$\bar r/r^2$ on returning to $z$. All other nonconstant entries scale
by factors with nonnegative power of $|r|$.

At $p=q=0$ the interpolant of $\bar z$ is $(\bar r/r)z$. To first
order in $q$, the root increments at $0,r,-r$ are $q/w,-q/(2w),-q/(2w)$.
Subtract the old interpolant applied to those increments from their
conjugates. The three differences are $a_0,-a_0/2,-a_0/2$ with
$a_0=\bar q/\bar w-(\bar r/r)q/w$. Their quadratic coefficient is
\begin{equation}\label{p10-eq:node-modes}
            -\frac{3a_0}{2w}=-\frac32\left(
                 \frac{\bar q}{|w|^2}-\frac{\bar r}{r}\frac{q}{w^2}
                              \right).
\end{equation}
The second term has angular mode $-3$ in $w$ and hence zero mean. In
either term subtracting the value of the test at zero costs $O(|q|)$
in the integral by H\"older continuity. Against the constant test one,
\eqref{p10-eq:node-modes} integrates to
$-6\pi\bar q\log(1/\ell)+O(|q|)$.

For completeness the Taylor remainders can be estimated before angular
integration. Writing $R=|r|$, the largest is bounded by
$C R^{-1}(|p|/R+|q|/R^3)^2$. Integration against $O(R^3\,dR)$ gives
\[
       O\bigl(|p|^2+|pq|/\ell+|q|^2/\ell^3\bigr)
                        =O(|p|^2+|q|).
\]
The $q$-linear terms of all other entries have integral $O(|q|)$ in
absolute value. The $p$-linear terms are bounded by $C|p|/|w|$ on
the disk, and deleting its inner part costs $O(|p|\ell^2)$. For
$\ell=|q|^{1/3}$ the logarithmic integral above is $-2\pi\bar Q+O(|q|)$.
For $|q|<|p|^3$ both $|q|\log(1/\ell)$ and $|Q|$ are uniformly
$o(|p|)$. Since $\ell^3\le |p|^3+|q|$, this proves, also on the axes,
that $M_{p,q}-M_{0,0}$ has $L^1$ norm $O(|p|+|Q|)$ and that its
pairing with any fixed H\"older test is its integral of its first-order
$p$-variation plus the point mass
\begin{equation}\label{p10-eq:node-mass}
              -2\pi\bar Q\,\delta_x
         \quad\text{in the }(\bar z,z^2)\text{ coefficient},
\end{equation}
up to $o(|p|+|Q|)$. Translating the branch center before this
calculation leaves the matrix on the smaller disk unchanged in form.
Variation of its smooth factors and of the data on the collars adds
their ordinary linear terms in $s,p$ and an error
$o(|s|+|p|+|Q|)$ in these pairings.

To pass to the obstruction map, pair with a central adjoint solution.
At the central parameter the conjugation matrix is bounded, sending
$\bar z$ to $(\bar r/r)z$ and $\bar z^2$ to $(\bar w/w)z^2$.
The adjoint solution, initially in $L^{a/(a-1)}$, thus lies in
$W^{1,k}$ for every finite $k$ by elliptic estimates and Sobolev
iteration. In particular it is H\"older continuous. The normalized
stabilized generator converges in $W^{1,a}$ to the central one.
In pairing the operator variation with the difference of these
generators the multiplier has $L^1$ norm $O(|s|+|p|+|Q|)$ and the
generator difference is $o(1)$ in $C^0$. The smooth variation of the
principal part on the collars gives the same little-oh estimate from
$W^{1,a}$ convergence. The central operator applied to the difference
pairs to zero. A varying stabilizing form contributes only its smooth
variation times its coefficient, which tends to zero. We may therefore
apply \eqref{p10-eq:node-mass} to the central generator in the actual
reduced equation. Its coefficient at the node is $B(x)v$.

\subsubsection{Identification of the residue and the local degree}

The distribution $a_1 z^2\delta_x$ represents in
\eqref{p10-eq:node-matching} a common nonzero real multiple of the matching
class of $a_1\in N_x$. To verify this, choose a local real-linear frame
$S(w)$ of $D$-homogeneous sections, with $S(0)=1$. Such a frame follows
by solving the local equation for the constant real basis vectors with
the Cauchy integral right inverse on a sufficiently small disk, and
renormalizing their values at zero. On the punctured disk the section
$z^2 S(w)a_1/w$ takes the value $S(w)a_1$ on the double and zero on the
identity. Multiply by a cutoff equal to one near zero. Its
distributional image is the ordinary cutoff forcing plus
$\pi a_1 z^2\delta_x$, since $\partial_{\bar w}(1/w)=\pi\delta_x$.
The ordinary forcing represents the matching class up to the convention
for the evaluation difference and projects to zero on $P_f$. The total
image pairs to zero with every adjoint by integration by parts, using
$1/w\in L^k$ for $k<2$ and the adjoint regularity just proved. This
identifies the residue in its own matching fiber, with the same real
scalar for every $a_1$.

On the nodal stratum $q_i=0$ the normalized component maps are
independent of the $p_i$. Their obstruction in $P_f$ has derivative
$W_f$ in $s$. The $i$th evaluation difference has derivative
$\pm v_i p_i$ in its moving mark, and may have a linear term in $s$.
These derivatives follow by stabilizing the normalized component
operators off the marks and varying the evaluation row in
\eqref{p10-eq:node-matching}. The comparison with the fixed-target
stabilization preserves the derivatives at the zero, since a continuous
change of trivialization only multiplies their linear part by its
central value.



\begin{proof}[Proof of Theorem~\ref{p10-thm:nodes}]
Choose the stabilizing forms away from all selected disks. Each disk
contributes \eqref{p10-eq:node-mass} with its own cutoff, so none of its
residue projects to another matching fiber. With $Q_i=q_i\log(1/|q_i|)$
extended by zero, the combined error is
\[
             o\left(|s|+\sum_{i\in I}(|p_i|+|Q_i|)\right)
\]
because there are finitely many disjoint disks. This holds on the
partially smoothed strata as well. After choosing a generator of $K$
and filtering each $N_{x_i}$ by $\CC v_i$, the linear part in these
coordinates is block triangular. Its diagonal blocks are $W_f$
evaluated on the generator, the complex-linear maps
$p_i\mapsto\pm v_i p_i$, and the maps
$Q_i\mapsto c_i\bar Q_i[B(x_i)v_i]$ to $N_{x_i}/\CC v_i$ for
nonzero real $c_i$. By \eqref{p10-eq:node-jet} the last maps are
antilinear isomorphisms of complex lines. The radial coordinate
change $q_i\mapsto Q_i$ preserves orientation near zero. The
little-oh estimate gives isolation and a homotopy on every sufficiently
small link to the invertible linear part.

Normalization gives the Fredholm determinant isomorphism
\[
  \det h_I^*D\simeq \det f^*D\otimes\det D\otimes
       \bigotimes_{i\in I}(\det_{\RR}N_{x_i})^* .
\]
It preserves the orientations obtained by deformation to the complex
normalization sequence. The evaluation fibers have complex orientation
and the one invertible identity component contributes $\epsilon$.
Replacing the generator of $K$ negates an even number $4j$ of real
matching coordinates. Each attaching block has positive sign and each
smoothing quotient negative sign. Thus the degree with the absolute
orientation upstairs is $(-1)^j\epsilon\eta_f$. On its link the
cokernel is the quotient by the image of the normalized obstruction
vector, hence pulls back the tangent bundle of the sphere of dimension
$2b+4j$. Its Euler number gives twice the local degree. Dividing by
the order two of the single deck group gives the asserted change.
Distinct subsets give distinct maps since the branch values in the
fixed target are distinct.

It remains to verify genericity. On a fixed double cover mark a simple
ramification point $r$ over $x$ and consider, on odd sections, the map
\[
          (D,\xi)\longmapsto(f^*D\xi,\partial_z\xi(r),B_D(x)) .
\]
It is a submersion at a nonzero solution. Indeed an annihilator of its
differential has an adjoint distribution solving the homogeneous adjoint
equation off $r$. On an unramified disk disjoint from $x$ where $\xi$ is
nonzero on one sheet, real zeroth-order variations prescribe arbitrary
forcing on that sheet and its negative on the other. Oddness and unique
continuation therefore force the distribution to be supported at $r$.
For any prescribed first holomorphic jet, choose a section whose equation
vanishes to higher order than the distribution there. This is obtained
by prescribing the pure holomorphic Taylor coefficients and recursively
solving the equation for coefficients involving $\bar z$, then cutting
off and antisymmetrizing. These variations kill the first-jet multiplier.
The adjoint distribution now solves the homogeneous equation across $r$,
so is smooth and zero. Free variation of $B_D(x)$ kills the remaining
multiplier. In Sobolev spaces of sufficiently high order the linearized
range is closed, since the section operator is Fredholm and the two
evaluations are finite dimensional.

At $v\ne0$, variation of $B$ surjects onto $N_x/\CC v$. Thus failure
of \eqref{p10-eq:node-jet} has real codimension two. The marked projectivized
odd incidence has index $-1$ over the operator space. Indeed, the odd
operator has real index $-2b$, since the invariant part has index zero
by Dolbeault descent. Varying the Hurwitz parameters adds $2b$ real
dimensions and projectivizing removes one. The branch mark has finite
choices. Adding a path parameter leaves the bad incidence of negative
expected dimension. Parametric transversality, followed by countable
intersection over covering genera and charts, gives the assertion
relative to rigid endpoints.
\end{proof}
\subsection{A ramification bound and compactness}\label{p10-sec:compactness}
Let $N$ for the moment be a Hermitian bundle of any rank over $C$.
Fix a unitary connection $\nabla$ and a metric on $C$ in its
conformal class. The graph equation for a lift $u$ of a finite
holomorphic map $f:Y\to C$ is
\[
  \bar\partial_{f^*\nabla}u + f^*P(\cdot,u)=0.
\]
Here $P$ is a smooth fiber-dependent $(0,1)$-form on a bounded
fiberwise convex tube. For a positive orthonormal pair $X,JX$ on $C$,
let $M$ bound the fiber Lipschitz constant of $P(\cdot,\cdot)(X)$,
uniformly in $X$. Put
\[
  \kappa=\bigl(\lambda_{\max}(iF_\nabla(X,JX))+2M^2\bigr)\,dA .
\]
We use the Hermitian product conjugate-linear in the first variable.
For a compact family of linear normal equations one may use a common
$M$, independent of the height of the graph.

\begin{lemma}\label{p10-lem:curvature}
Suppose $(f,u)$ is generically one-to-one on its smooth
normalization, with distinct component images. Let $m=\deg f$,
let $R$ be its total ramification, and suppose $n$ disjoint pairs
of distinct normalization points are identified as nodes at which
the graph matches. Then
\[
 R+2n\le \frac{m(m-1)}{2\pi}\int_C\kappa .
\]
In particular a connected graph with these nodes and no contracted
components has arithmetic genus at most
$1+m(g-1)+\frac{m(m-1)}{4\pi}\int_C\kappa$.
\end{lemma}
\begin{proof}
Normalize the closure of the off-diagonal pairs in $Y\times_C Y$ over
regular values of $f$. The resulting smooth, possibly disconnected
curve $Z$ has projections $p_1,p_2$ and a map $b=fp_i$ of degree
$m(m-1)$ to $C$. The section $s=p_1^*u-p_2^*u$ is nonzero somewhere on
every component. Subtracting the two equations gives
$|(\bar\partial_{b^*\nabla}s)(T)|\le M|s|\,|db(T)|$ for each unit
tangent $T$. Integrating the vertical derivative of $P$ along
the segment between the two heights writes the difference as a
smooth real-linear zeroth-order operator applied to $s$. Its zeros
are isolated and have finite order by the similarity principle.
In a smooth unitary frame the leading Taylor term of order $a$ is
$z^a v$ with $v\ne0$, since the principal equation on that term is
$\bar\partial=0$.

Away from the zeros, set $e=s/|s|$ and
$\alpha=\operatorname{Im}\langle e,b^*\nabla e\rangle$.
Project the derivative perpendicular to $\CC e$ and denote the
result by $\nabla^\perp e$. Differentiation in a positive unit
frame $T,jT$ gives
\begin{align*}
 d\alpha(T,jT)={}&\operatorname{Im}
   \langle e,(b^*F_\nabla)(T,jT)e\rangle \\
  &+2\bigl(|(\nabla^\perp e)^{1,0}(T)|^2-
                            |(\nabla^\perp e)^{0,1}(T)|^2\bigr).
\end{align*}
The components parallel to $e$ have purely imaginary coefficients and
drop out of the exterior product. The bound on $\bar\partial s$ gives
$|(\nabla^\perp e)^{0,1}(T)|\le M|db(T)|$, whence
$d\alpha\ge-b^*\kappa$. On a positively oriented small circle about
a zero of order $a$, the integral of $\alpha$ tends to $2\pi a$.
Stokes' theorem on the complement of the disks bounds the sum of the
orders by $(2\pi)^{-1}\int_Z b^*\kappa$.
At a point of index $k$ there are $k-1$ off-diagonal branches over
the repeated point, one for each nontrivial $k$th root of unity.
Each forces a zero. Each identified pair forces a zero in each
order of the pair, at points distinct from the ramification pairs.
This proves the inequality. Riemann--Hurwitz on the components
and the arithmetic genus formula for a nodal curve give the last claim.
\end{proof}

\begin{proposition}\label{p10-prop:compact}
Over a generic compact path as in Theorem~\ref{p10-thm:chambers}, the
projectivized incidence of smooth connected cubic covers with kernel
is compact and finite, with the genera of all covers included.
\end{proposition}
\begin{proof}
Proposition~\ref{p10-prop:incidence} gives discreteness. For a sequence of
kernel lines choose generators of fixed small maximum norm and mark
points attaining that norm. Their embedded graphs lie in a compact
subtube. The family of graph almost complex structures agrees there
with a smooth family on a compact completion, tamed by a common form
as in the proof of Proposition~\ref{p10-prop:smoothjump}. Projection
determines the homology class and hence gives a uniform bound on
area. Lemma~\ref{p10-lem:curvature} bounds the genus, so a subsequence
has fixed genus. Stable-map compactness gives a limit whose marked
point has nonzero height. Forget that mark and stabilize.

A nonconstant component projects in positive degree, since the
portion of a single fiber in the tube lies in a complex vector space.
Invertibility on $C$ excludes a degree-one component of nonzero
height. If a component has degree three and nonzero height, its
parametrization is simple, since the alternative factors through
a nonzero graph of degree one. This cubic normalization is embedded
by Proposition~\ref{p10-prop:incidence}. Otherwise the positive components
are a nonzero double graph and one copy of the zero section,
both with multiplicity one. The double generator is odd. It can vanish
only at ramification points, since a zero over a regular value would
identify the two sheets of its embedded graph. Its derivative at a
ramification point gives a vertical tangent, distinct from the tangent
to the zero section.

For a maximal stable contracted subcurve in a limit of smooth
holomorphic curves, the attaching branches include a critical branch,
two branches with coincident tangent lines, or at least three
branches \cite[Theorem~1.1]{ekholm}. This theorem applies to the
varying structure by extending the parameter interval to a complex
disk, smoothly extending the family in its real coordinate, and
using the fixed block structure $J_t\oplus i$ on the product.
The curves are constant in the parameter direction. Each candidate
limit above has at most two branches through any image point, and
they are immersed with distinct tangent lines when both occur.
Consequently there is no contracted subcurve. Neither embedded
normalization admits a self-node. A single cubic limit therefore
converges in the smooth-cover chart and retains its kernel line.
In the two-component case connectedness requires direct joins at a
nonempty subset of the ramification points. Theorem~\ref{p10-thm:nodes}
isolates the projective kernel there even in the full smoothing
chart. Such a point cannot be approached by smooth cubic solutions.
Rigidity excludes endpoint limits. This proves compactness.
\end{proof}

For degree two a normalized nonzero graph has only the embedded
degree-two possibility in the preceding proof. The same argument
gives finiteness, also proved in \cite[Proposition~3.73]{dw}.
The contracted components excluded in Proposition~\ref{p10-prop:compact}
do enter the Euler class on the full stable-map space.
\subsection{Contracted components at a reducible cubic cover}\label{p10-sec:descendants}

Let $f:Y\to C$ be the double cover at a crossing as in
Theorem~\ref{p10-thm:nodes}, and take the crossing parameter to be $0$.
Write $r_1,\ldots,r_b$ for its ramification points, $x_i=f(r_i)$,
and $\xi$ for a nonzero odd kernel section. For a nonempty set
$I\subset\{1,\ldots,b\}$, let $h_I:X_I\to C$ be the finite map
obtained from $Y\amalg C$ by identifying $r_i$ with $x_i$ for
$i\in I$. Put
\[
 H=g(Y)+g(C)+|I|-1,\qquad d_I=(-1)^{|I|}\epsilon\eta_f .
\]
In arithmetic genus $k\ge H$, let $G_{I,k}$ consist of the stable
maps obtained from $h_I$ by attaching connected contracted curves
at smooth points and by inserting connected contracted curves between
the branches of its nodes. We include all stable configurations in
these operations. In particular an inserted curve may have rational
components. We regard $G_{I,k}$ as a locus at $t=0$ in the product
of the ordinary fixed-target moduli stack with the parameter interval.

\begin{theorem}\label{p10-thm:descendants}
Suppose the hypotheses of Theorem~\ref{p10-thm:nodes} hold with its
projected jet condition at every ramification point. Then $G_{I,k}$
is compact and isolated in the locus of stable maps with nonzero
pullback kernel.
Its contribution to the change of the absolute normal Euler number in
arithmetic genus $k$ is
\begin{equation}\label{p10-eq:desc-contribution}
 d_I\,[v^{2k-2}]\bigl(2\sin(v/2)\bigr)^{2H-2}.
\end{equation}
Thus the contribution of the specified set of joining nodes to the
connected series is $d_I x^{H-1}$, where $x=(2\sin(v/2))^2$.
\end{theorem}

We displace the nodal kernel to a signed collection of embedded
graphs on smooth domains. Near each such graph the Euler integral is
a constrained virtual count in the total space. Cutting its parameter
interval at one regular value leaves the ordinary degree-one series.


\subsubsection{The compact locus and the choice of sheets}
A presentation by proper gluing maps makes explicit the allowable
specializations at an attachment. Choose nonnegative integers
$e_Y,e_C,e_i$ with $e_Y+e_C+\sum_{i\in I}e_i=k-H$.
Take the two evaluation fibers
\[
 \overline{\mathcal M}_{g(Y)+e_Y,|I|}(Y,1)_{(r_i)_{i\in I}},
 \qquad
 \overline{\mathcal M}_{g(C)+e_C,|I|}(C,1)_{(x_i)_{i\in I}} .
\]
At $i$, join the corresponding marks directly if $e_i=0$ and
insert a member of $\overline{\mathcal M}_{e_i,2}$ if $e_i>0$.
Compose the maps to $Y$ with $f$. The union of the images over these
finitely many choices is $G_{I,k}$. Indeed a degree-one stable map
to a smooth curve has one nonconstant normalization, mapping
isomorphically, and every connected contracted subcurve maximal for
inclusion meets this normalization just once. Its markings may lie
on contracted components. This accounts also for tails whose attaching
points come together or specialize to a joining point. Conversely,
cutting out all maximal contracted subcurves that join the two
positive components, and marking the ends of each direct joining node,
recovers such gluing data.

In the lift to $N$, the positive components are the graph of $\xi$
over $Y$ and the zero section over $C$. They are embedded and meet
precisely over the $x_i$, with distinct complex tangent lines.
Every contracted connected subcurve has a single image point. Thus
its external positive branches number at most two, and two can occur
only on opposite components. The direct nodes and the contracted
subcurves with two ends recover $I$ intrinsically. The proper gluing
images for distinct $I$ are therefore disjoint. Restriction to the
positive components identifies the kernel throughout $G_{I,k}$ with
$\RR\xi$, since the base operator is invertible and on contracted
subcurves the sections are constants with prescribed attaching values.

Choose concentric closed disks $U\subset\operatorname{int}U'$ of
regular values disjoint from the $x_i$. Over $U'$ the three heights
$0,\xi,-\xi$ are separated. Even if a contracted curve lies over a
point of $U$, a nearby stable map has three pieces over $U$, each of
total degree one. In fact at a central map the three preimages of $U'$
with their contracted components have disjoint neighborhoods. The
collapse maps of converging domains keep the preimage of $U$ in
these neighborhoods, and evaluation at a regular value away from the
limiting nodal images gives degree one in each. Order the pieces with
the identity sheet first. These two orderings over the central locus
extend to a finite étale cover of a neighborhood, by taking all sheet
orderings locally and continuing the two selected ones. Its stack
degree is two. The deck involution of $f$ exchanges the orderings
and fixes each $r_i$. This removes the double-cover symmetry at the
finite map and retains the isotropy of the contracted curves. We work
on this cover until the last step of the proof.

We recall the local class computing the Euler change. In a finite
stabilization $S:E\to F$ of the normal family, projectivize $E$
over $\RR$ and restrict $S$ to its tautological line. The pushforward
of the resulting relative Euler class is supported on the kernel
locus. Near $G_{I,k}$ we can exclude kernels of dimension two by
upper semicontinuity. In the space of matrices the normal space at a
matrix of kernel dimension one is $\operatorname{Hom}(K,P)$ for
its kernel $K$ and cokernel $P$. The pushforward gives the Thom
generator there, since the projective vertical differential cancels
the image of the matrix. On the normal sphere the oriented cokernel
is its tangent bundle after choosing a generator of $K$. The sphere
has even dimension, and its Euler integral is two
\cite[Lemma~6.10]{bai}. Thus the connecting class of the cokernel
Euler class off the kernel locus equals twice the projective class
by the Thom isomorphism in the matrix space and pullback. Its pairing
with the fixed-target virtual class times the oriented interval is
the terminal Euler number minus the initial one. Both this identity
and its restriction by excision are preserved by common stabilization.

\subsubsection{Displacement on the finite-map charts}
On one labelled double sheet let $\ell_a(u)$ be the integral of a
smooth real covector applied to $u$, against the pullback by $a$ of a
cutoff area form supported in $\operatorname{int}U$. Choose the test
so that $\ell_a(\xi)\ne0$. It gives the same normalization of the
kernel through every attachment on that sheet, since pullback area
vanishes on contracted components. Let $j_a(e)$ denote a linear
combination of triples of smooth $N$-valued $(0,1)$-forms supported
in $\operatorname{int}U$, pulled back separately on the corresponding
labelled pieces. For the variable stable map $a$ put $L_{a,t}=a^*D_t$
and consider the finite-rank perturbation
\begin{equation}\label{p10-eq:desc-perturb}
 T^e_{a,t}u=L_{a,t}u-\ell_a(u)j_a(e).
\end{equation}
For continuity of the reduction at contracted curves we explain how
to use the nodal stabilization of \cite[Appendix~B]{bai}. Add a
scalar $z$ with $L_{a,t}u=zj_a(e)$ and $z=\ell_a(u)$. Fix a central
stable map. Enclose the distinct nodal values meeting the supports of
the forcings in small disjoint disks inside $U'$. On each disk and
each labelled piece choose a smooth primitive $b_i$ solving
$D_t b_i=j_i(e)$ on the base. These primitives can be chosen
smoothly in $t$ and linearly in $e$: on a small disk the Cauchy
right inverse for $\bar\partial$ followed by the zeroth-order term
has norm less than one after rescaling. With cutoffs equal to one on
smaller disks about the nodal values, pull the primitives back on
their respective labelled pieces and extend by zero. Write $b_a$ for
their sum. The substitution $u=u'+zb_a$ changes the forcing to
$j_a(e)-L_{a,t}b_a$ and the scalar equation to
$z=\ell_a(u')+z\ell_a(b_a)$. The new forcing vanishes on
neighborhoods of all the nodes. This uses the pullback identity on
maps to $C$, including the constant components.

On shrinking the chart, the remaining forms vary smoothly on compact
thick subsets away from the necks in the sense of convergence of
stable maps. In a domain trivialization they converge there in every
smooth norm; pullback by $a$ involves only its first derivative.
One may freeze these forms at the center in the nodal stabilization:
the error tends to zero in operator norm and is absorbed by the
uniformly bounded stabilized inverse
\cite[Propositions~B.7.9 and~B.10.4]{pardonatlas}. At the center a minimal cokernel
supplement $B$ gives the invertible augmented operator
$(v,w)\mapsto(L_{a,t}v-Bw,\ell_a(v))$ since $\ell_a$ is nonzero on
the kernel line. The area test on shrinking base disks has size at
most a constant times their area for bounded sections, by the
degree bound. This proves continuity also of the scalar row. On
overlaps the substitutions are translations by the difference of
the two sums of primitives times $z$, with the induced scalar row,
and hence give the same projective class. Near a central normalized line, its stabilized sections
converge on thick regions and are uniformly bounded on the domains.
Thus the area test is continuous, and the normalization $\ell_a(u)=1$
is available near the central projective kernel. On this slice the
equation becomes
\begin{equation}\label{p10-eq:desc-affine}
                  L_{a,t}u=j_a(e),\qquad \ell_a(u)=1 .
\end{equation}
The first equality is a local graph equation. In the three disjoint
fiber neighborhoods over the support of the forms, multiply the
corresponding forcing by a smooth cutoff equal to one near that
sheet. The resulting affine term in the horizontal-to-vertical
block defines an almost complex structure on $N$ with holomorphic
projection and standard fibers \cite[Lemma~5.9]{bai}. On the
normalized graphs it gives the indicated equation even when a
contracted curve lies over the support. Cutting off outside a larger
disk bundle extends the structure over $\PP(N\oplus\CC)$.
The fiberwise Fubini--Study form extended by a connection, plus a
sufficiently large pullback of an area form on $C$, tames the family.

Fix $k$. Use the same sheet data on the finite-map charts for all
$h_J$ of genus at most $k$, $J\ne\varnothing$. By
Theorem~\ref{p10-thm:nodes} each has an isolating closed neighborhood in
its own genus. Independent forms on the labelled sheets span the
cokernel at every one of their centers. To see this, present the
operator on the normalization with one evaluation-difference row
at each join. A dual annihilator of the sheet forms vanishes on
a disk of each positive component. The adjoint equation and unique
continuation give vanishing off the joins. The equation-space
multipliers lie in the dual of $L^p$ for $p>2$, so none is supported
at an isolated point. Sections with independently prescribed endpoint
values then kill the remaining evaluation multipliers. Hence one
finite-dimensional space of triples suffices for all these centers.

In a minimal stabilization the normalized obstruction section takes
values in this space of representatives, and \eqref{p10-eq:desc-affine}
translates the section by the constant parameter $e$. If more forms
are used, the derivative in the forcing parameters is still
surjective. On each stratum of fixed nodal and branching type the
section is smooth by normalization with the matching equations.
Persistence of a node has positive complex codimension by the local
smoothing parameters of Theorem~\ref{p10-thm:nodes}. On smooth sources
collision of branch values also has positive complex codimension by
the Hurwitz dimension formula. The unrestricted normalized section
has expected real dimension zero at fixed $e$, with the path
parameter included. Regular values therefore avoid these strata and
give isolated zeros on smooth simply branched coverings. Compactness
of the isolating charts makes their number finite.

We may arrange also that their graphs are embedded. On a smooth
source the vertical derivative at a simple ramification point takes
values in a complex two-plane. Equality of heights at two distinct
points over a regular value imposes four real conditions with two
real parameters for the value. At a branch value the marks are
discrete. The corresponding marked evaluation and first-derivative
conditions are transverse under the sheet forcings. Indeed an
annihilating distribution vanishes on the open test sheets away
from the marks by the independent variations. The adjoint equation
there eliminates any multiple of the normalization density by testing
where that density is nonzero. Unique continuation leaves a
distribution supported at the marks. In a local trivialization one
can prescribe independent values and complex first derivatives at
the marks with all jets of the linearized equation vanishing up to
any specified finite order. This follows by solving recursively for
the derivatives involving $\bar z$ from the Cauchy--Riemann symbol.
Testing on these sections removes the point multipliers. One first
obtains immersion, which excludes coincidences in a neighborhood
of the diagonal. On the compact remaining marked pairs, finitely many
additional forms give the required transversality near the finitely
many regular solutions. Arbitrarily small regular choices in the
smooth space of triples result, each in a finite-dimensional
perturbation. Write $\sigma_\Gamma$ for the oriented projective
degree of a resulting smooth solution in the chart of $h_I$.
Conservation of the degree of the normalized obstruction section in
that chart gives
\begin{equation}\label{p10-eq:desc-sum}
                        \sum_\Gamma\sigma_\Gamma=d_I
\end{equation}
by the oriented local calculation of Theorem~\ref{p10-thm:nodes}.

\subsubsection{Removing the contracted components}
Consider a sequence of projective zeros tending to $G_{I,k}$ with
$e\to0$, allowing also $e=0$. Use the normalized graph equation and
take stable limits of the nonconstant normalizations separately,
first retaining their nodal flags. Their limiting nonconstant images
together are the graph of $\xi$ and the zero section, each with
multiplicity one. This follows by gluing the marked limits and
forgetting the flags. Each nonconstant normalization retains positive
projection degree, since every compact curve in an affine fiber is
constant. Thus there are either one or two such normalizations.

In the unmarked stable limit of each normalization there are no
contracted components. In fact ghost bubble censorship requires a
vanishing derivative at a positive attaching branch, equal complex
tangent lines at a pair of attaching branches, or at least three
positive attaching branches at one point
\cite[Theorem~1.1]{ekholm}. The two central embeddings exclude
all three possibilities. To use the fixed-structure theorem in a
smooth finite-dimensional family, extend the parameters locally to
an open subset of $\CC^m$ and put the block-diagonal almost complex
structure on the product. The graphs on constant parameter slices
are holomorphic for this single structure, and the attaching
directions are unchanged. If a sequence uses successively larger
spaces of smooth forcings, choose a subsequence whose first $n$
seminorms differ from the limit by at most $2^{-n^2}$ at step $n$.
Disjoint smooth bumps of width $2^{-n-3}$ centered at $2^{-n}$
interpolate it to a smooth arc flat at zero. The same product argument
then applies.

If there are two positive normalizations, their individual unmarked
limits are smooth embedded copies of $Y$ and $C$. An identification
on either limit would give a self-intersection of its embedding.
The normalizations thus converge smoothly and remain embeddings.
Near a point of intersection of the two limiting images, project
the difference of their two local parametrizations onto the real
four-plane spanned by the two tangent lines. Its derivative in the
two source variables is invertible. On smaller neighborhoods the
inverse function theorem gives at most one intersection for
each nearby pair, and there are none away from these neighborhoods.
Consequently every maximal contracted subcurve has one end on
a positive normalization or one on each. Delete the one-ended
subcurves, replace the two-ended ones by nodes, and retain the
direct nodes. Matching and the thick integral show that the finite
nodal map so obtained satisfies the same normalized equation.
After subselection its joining sites approach a fixed subset $J$.
The flags for different connections have distinct limits on either
component. Marked stabilization and properness of the one- and
two-pointed contracted factors recover a limit in $G_{J,k}$, even
when other flags coalesce with a joining flag. Hence $J=I$ and the
finite zeros converge in the chart of $h_I$. At $e=0$ they are
$h_I$ itself, by Theorem~\ref{p10-thm:nodes}. The regular displacements
just chosen avoid all finite nodal zeros there.

With one positive normalization, its unmarked limit is a connected
finite union of the two central embeddings with direct nodes at some
nonempty subset $J$. Its normalized equation is unchanged on deleting
the other components, which are constant. For $e=0$, smooth finite
zeros approaching $h_J$ contradict Theorem~\ref{p10-thm:nodes}. For the
chosen regular displacements they are smooth embedded graphs from
the chart of $h_J$. Embeddedness excludes self-nodes and forces each
contracted connected subcurve maximal for inclusion to attach once.
Mark these attachments before passing to the limit. Forgetting the
marks on the positive limit recovers $X_J$. Restoring them can insert
a rational tree over a smooth point or between the two branches of an
existing node. Proper limits of the one-pointed tails attach to these
trees. In particular flags approaching two distinct smooth points
above an unjoined graph intersection stay on distinct fibers of the
stabilization. The limit belongs to $G_{J,k}$, so again $J=I$.

The zero-displacement argument proves isolation. Properness in fixed
genus and closedness of the kernel locus give a relatively compact
neighborhood with kernel-free boundary, which persists through
sufficiently small perturbations. The two cases show that within it
the displaced zeros are precisely the stable degree-one expansions
of the smooth graphs from the chart of $h_I$. For the reverse inclusion,
every such expansion solves \eqref{p10-eq:desc-affine} by constant
extension. If expansions escaped the neighborhood as $e\to0$, their
marked stable limits would lie in $G_{I,k}$ by the same gluing
argument. This also separates the compact families over the individual
displaced graphs.

\subsubsection{The virtual class of the normalized graph equation}
Fix one displaced graph $\Gamma$, at $t_0$, and write $a_0$ for its
projection. The homogeneous operator $L_{a_0,t_0}$ is injective.
Indeed $a_0$ has smooth source in the isolating finite-map chart of
the original nodal zero. Its pullback to every degree-one expansion
is still injective, by restriction to the positive normalization
and matching constants on the other components. On this compact
family the cokernel has real rank
\[
                   r=4\bigl(k-1-3(g(C)-1)\bigr),
\]
the real virtual dimension of $\overline{\mathcal M}_k(C,3)$.
We compare the projective Euler pairing for \eqref{p10-eq:desc-perturb}
with the virtual degree of ordinary maps to the almost complex
total space subject to the scalar equation in \eqref{p10-eq:desc-affine}.

Take an equivariant chart for maps to $C$, with stabilizing markings,
incidence slices, horizontal perturbation forms, and obstruction
section $s_C$. Use the same markings for maps to the total space,
taking inverse images of the target slices. Lift the horizontal forms
using a complex-linear splitting of projection for the graph
structure. On a horizontally thickened map the vertical equation
is still affine in $u$. Indeed in connection coordinates the
off-diagonal term of the almost complex structure is affine in height
where a sheet forcing acts, and linear elsewhere near the graph.
The horizontal lift can be chosen affine in height there by projecting
a connection lift to its complex-linear part. Taking the
$(0,1)$-part on the source then gives an affine vertical equation
with the usual Dolbeault symbol, restricting to $L_{a,t}u-j_a(e)$ on
$s_C=0$.

Choose vertical forms $B w$, $w\in W$, inducing an isomorphism from
$W$ to the cokernel of $L_{a,t}$ at the chart center. Forms on the
smooth interiors of every component, including contracted ones, are
available. They span by density off the nodes in the equation spaces
for sections with matching values. Choose equivariant representatives
by averaging and extend them over the stabilized universal domain.
A contracted component can require a vertical stabilizer even though
the forcing $j_a(e)$ vanishes there. In these conventions the
vertically stabilized equation on $s_C=0$ is
\begin{equation}\label{p10-eq:desc-stabilized}
                      L_{a,t}u-Bw=j_a(e).
\end{equation}
The derivative $(L_{a,t},-B)$ in $(u,w)$ is an isomorphism at the
center. In the simultaneous thickening the first row is the regular
stabilized horizontal equation, independent of the height. The
second has this isomorphism as its vertical derivative. Projection
therefore identifies the ambient thickening locally with the
horizontal one, with $(u,w)$ uniquely determined on it. This holds
over a nodal fiber by the regular perturbed-map chart and inverse
function theorem relative to the common base of stabilized domains
\cite[Theorem~3.5 and Lemma~5.10]{swam}. The ambient obstruction
section on the common chart is $(s_C,w)$.

On the horizontal zero set introduce $z=\ell_a(v)$ for a homogeneous
section $v$. Eliminating the variables in
$L_{a,t}v-Bw'=zj_a(e)$ by \eqref{p10-eq:desc-stabilized} gives
$v=zu$, $w'=zw$. At a projective zero $w'=0$ implies $z\ne0$ by
injectivity of $L_{a,t}$. Thus in the chart $z=1$ the residual
equations are exactly
\begin{equation}\label{p10-eq:desc-rows}
                           w=0,\qquad \ell_a(u)-1=0 .
\end{equation}
Capping the horizontal virtual class with the localized projective
Euler class amounts on this common thickening to multiplying the
Thom classes of $s_C$ and of \eqref{p10-eq:desc-rows}. Their product is
the class of the constrained ambient equation.

For compatibility, on overlaps take the union of stabilizing marks
and the sum of the horizontal and vertical perturbation spaces,
transporting earlier forms by the forgetful maps of marked domains.
The lifted slices and projection commute with these forgetful maps.
An enlarged vertical space adds free variables in
\eqref{p10-eq:desc-stabilized} paired with identical additional
obstruction coordinates. Switching the minimal space $W$ changes
cokernel coordinates and adds a multiple of $w$ to the scalar row.
A change of horizontal lift adds terms vanishing on $s_C=0$.
The resulting stabilizations and triangular row operations preserve
the localized Thom product, since throughout the homotopies the
joint zeros stay in the compact interior. The same sums refine triple
intersections, equivariantly throughout, as in the ordinary
Gromov--Witten implicit atlas \cite[\S9.2, Theorem~9.2.10]{pardonatlas}.
The determinant orientation under projection is the product of the
horizontal and vertical orientations, by the triangular
linearization and homotopy to its complex symbol. Moving the scalar
row past the horizontal virtual dimension or the vertical rank $r$
has even parity. Its orientation convention can thus be fixed at
the genus-$H$ crossing and used in every arithmetic genus.

\subsubsection{Restriction to one parameter value}
At $\Gamma$ let $A$ be the fixed-time normal linearization of the
ordinary ambient equation. Its real index is zero since the ambient
first Chern number on the graph is zero. The derivative with time and
the integral constraint, in unparametrized normal coordinates, is
\begin{equation}\label{p10-eq:desc-augmented}
 (\zeta,\tau)\longmapsto
          (A\zeta+b_0\tau,\lambda\zeta).
\end{equation}
Here $b_0$ is the parameter derivative and $\lambda$ the differential
of the time-independent integral test. This operator is invertible.
Indeed the horizontal equation at a smooth finite cover is regular
with variable domain, as its deformation complex reduces to the
torsion cokernel of the derivative of the cover. Eliminating that
equation and the vertical complement leaves the transverse normalized
equation that defined $\sigma_\Gamma$. Hence $\dim_{\RR}\ker A\le1$.

If $A$ is singular, choose an unramified subdisk of a labelled sheet
on which generators of its kernel and adjoint cokernel are nonzero.
Such a subdisk exists by unique continuation. If $u_\Gamma(w)$
is its height, add to the affine equation a small term
\[
                     \delta B_1(w)\bigl(u-u_\Gamma(w)\bigr)
\]
with support in this subdisk and a fiber cutoff equal to one near
the sheet. On an unramified graph the vertical tangent projects
isomorphically onto the normal quotient. We can therefore prescribe
the zeroth-order normal variation so that its pairing between the
two generators is nonzero. The one-dimensional Fredholm reduction
makes $A$ invertible for sufficiently small $\delta\ne0$, leaves
$\Gamma$ a solution, and preserves invertibility of
\eqref{p10-eq:desc-augmented}. Near the graphs this is the homogeneous
deformation \eqref{p10-eq:desc-perturb} with $L$ and $j$ replaced by
$L+\delta B_1$ and $j+\delta B_1u_\Gamma$ on the chosen sheet.
By compactness and continuous stabilization its size can be chosen
to preserve vertical injectivity near the expansions and projective
nonvanishing on the boundary of their isolated neighborhood, both in
genus $k$ and in genus $H$. Thus it preserves the localized degrees.

To see that no part of this degree lies away from the retained graph
after the perturbation, consider a sequence of other constrained
solutions as $\delta\to0$ in the original isolated neighborhood.
Its compact limit is an expansion of $\Gamma$. In a tube retracting
to that multiplicity-one curve only one nonconstant normalization
can occur, since a positive component projects with degree a
positive multiple of three. The unmarked normalization then converges
smoothly by ghost bubble censorship at the embedded limit. The
implicit function theorem for the augmented operator
\eqref{p10-eq:desc-augmented}, with $\delta$ as parameter, gives the
unique nearby constrained curve, namely the retained $\Gamma$
at $t_0$. All remaining components attach to it at one point each.

At the regular endpoint the nearby ordinary embedded graphs form
a smooth family $\Gamma_t$. In a small tube retracting to
$\Gamma$, the projection of every positive class represented by a
holomorphic component has degree a positive multiple of three.
A map of total projection degree three has one nonconstant component,
of multiplicity one. At fixed genus, a sequence of its unmarked
normalizations approaching $\Gamma$ has smooth limit. Indeed the
limiting positive map is the embedding with multiplicity one, and
ghost bubble censorship excludes any contracted component of the
unmarked limit. Smooth convergence places the normalization in the
regular chart of $\Gamma_t$.
Compare also the embedded-cycle compactness and openness results
of \cite[Propositions~2.21 and~2.29]{diw}. Embeddedness then implies
that every other component belongs to a one-ended contracted subcurve.
Consequently, after shrinking the tube and parameter interval, the
entire fixed-genus ordinary locus near the compact family consists of
degree-one stable maps to $\Gamma_t$.

On this locus the integral constraint equals the function of time
$\gamma(t)=\ell(\Gamma_t)-1$, since every attaching subcurve has
zero pullback area. Invertibility of \eqref{p10-eq:desc-augmented} gives
$\gamma'(t_0)\ne0$ by applying its scalar row to the tangent of
the implicit family. Interpolating between $\ell_a(u)-1$ and
$\gamma(t)$ on the common thickening preserves the localized product,
as the two functions agree wherever $s_C=w=0$. The real Thom class
of $\gamma$ cuts the parametrized ambient virtual class down to the
ordinary fixed-time class with sign $\operatorname{sgn}\gamma'(t_0)$.

Write $D_\Gamma$ for the regular normal operator on $\Gamma$.
Its normal bundle has degree $2H-2$, since
$c_1(TC)+c_1(N)=0$. Every degree-one pullback of $D_\Gamma$ is
injective by restriction and matching. Zinger's comparison therefore
expresses the local degree as the normal Euler pairing on the ordinary
degree-one stable-map class of $\Gamma$ \cite[Theorem~1.2]{zinger}.
To record its orientation, deform $D_\Gamma$ on the fixed bundle.
For each degree-one map the pullback and base operators have the same
kernel throughout the deformation. On cokernels restriction to the
positive normalization is a left inverse to pullback of forms.
Hence the quotient of their two stabilized Fredholm complexes is
injective with vector-bundle cokernel throughout
\cite[Appendix~B]{bai}. Its determinant is the quotient determinant.
At the regular endpoint its Euler class gives the absolute pairing
multiplied by the base determinant sign $\epsilon(D_\Gamma)$.
At the Dolbeault operator of a sum of two non-effective line bundles
of degree $H-1$ the orientation is complex. Pandharipande's
degree-one Hodge integral \cite[Theorem~2]{pandhodge} gives its
generating series $\bigl(2\sin(v/2)\bigr)^{2H-2}$. The oriented
fixed-time series at the regular endpoint is therefore
$\epsilon(D_\Gamma)x^{H-1}$.

\begin{proof}[Completion of the proof of Theorem~\ref{p10-thm:descendants}]
The product of the time-slice sign and the fixed-time determinant
sign is $\sigma_\Gamma$, as follows by evaluating their product in
genus $H$ at the transverse normalized graph. The scalar orientation
and quotient determinant comparison above give the same product in
genus $k$. Thus this graph contributes
$\sigma_\Gamma[v^{2k-2}]x^{H-1}$ to the projective pairing upstairs.
Summing by \eqref{p10-eq:desc-sum}, multiplying by two for the full
Euler change, and dividing by the stack degree two of the sheet
labelling gives \eqref{p10-eq:desc-contribution}. All isolating
neighborhoods and perturbation sizes were chosen at fixed genus.
In particular the formula includes tails and bridges of every
arithmetic genus allowed by $k$.
\end{proof}
